
\documentclass[11pt]{article}
 
\usepackage[margin=1in]{geometry}
\usepackage{amsmath,amssymb,amsthm}
\usepackage{booktabs}
\usepackage{graphicx}
\usepackage{xcolor}
\usepackage{hyperref}
\usepackage{enumitem}
\hypersetup{colorlinks=true, linkcolor=blue, citecolor=blue, urlcolor=blue}
 
\newcommand{\Hid}{\mathrm{Hid}}
\newcommand{\Acc}{\mathrm{Acc}}
\newcommand{\Rc}{\mathcal{R}}
\newcommand{\Mrmu}{M_{\mathrm{RMU}}}
\newcommand{\Mbase}{M_{\mathrm{base}}}

\title{lab notebook week 2}
\author{Alessa Carbo}
\date{June 2026}

\begin{document}

\maketitle

\section{Introduction}
This week I'm running my probing and recovery dynamic experiments. For Monday June 8th, we've started with RMU and WMDP-bio.  We operationalize this
distinction through a recovery-budget quantity, $\Hid_\tau$, and instantiate its
\emph{representational} axis: we ask whether the suppressed answer to a WMDP-bio
multiple-choice question is decodable from the frozen residual stream of an RMU-unlearned
model (\texttt{cais/Zephyr\_RMU}) by a probe, even though the model's behavioral accuracy
is floored. Across a grid spanning two probe capacities (linear, MLP) and four
token-aggregation schemes (last-token, full-sequence mean, option-span pool,
learned-attention), at all 33 layers, no probe recovers the answer from the RMU model,
while the identical probes recover it from the base model in every cell.And we have a selectiv null because the same pipeline recovers general biology (MMLU-bio, not in the forget set)
from the RMU model as well as from base, ruling out collateral damage/buggy implemention or/and generic
probe-blindness. We can see that RMU's suppression is likely deep and targeted on the
representational-recovery axis at the probe budgets tested,though this is in onlu one possible axis/budget.


\section{Formal setup}
 
\subsection{The recovery-budget quantity}
Let $M$ be an unlearned model and $T$ a target task. Let $\Rc_c$ denote a
\emph{recovery procedure} parameterized by a budget $c$, mapping a model to a recovered
predictor for $T$. Following the project's formalization, define the \emph{hiddenness} of
the capability at threshold $\tau$ as the minimum budget at which recovery clears $\tau$:
\begin{equation}
\Hid_\tau(M, T) \;=\; \min\Bigl\{\, c \;:\; \Acc\bigl(\Rc_c(M), T\bigr) \ge \tau \,\Bigr\},
\label{eq:hid}
\end{equation}
with $\Hid_\tau(M,T) = \infty$ if no budget in the considered class reaches $\tau$.
Small $\Hid_\tau$ means the capability is shallowly hidden (a cheap recovery suffices);
large or unreached $\Hid_\tau$ means it is deeply removed or inaccessible to the recovery
class. The budget $c$ is not a single scalar: the proposal distinguishes at least a
\emph{representational} axis (the computational class of a probe reading frozen
activations), an \emph{optimization} axis (fine-tuning steps / weight movement), and an
\emph{in-context} axis (number of ICL demonstrations). We report the representational axis only for the moment.


\subsection{The representational recovery operator}
Fix a model $M$ with $L+1$ residual-stream sites (embeddings plus $L$ transformer layers),
hidden dimension $d$. For an input $x$ let $h^{(\ell)}_t(M,x)\in\mathbb{R}^d$ be the
residual-stream activation at layer $\ell$ and token position $t$. A representational
recovery procedure is a choice of (i) layer $\ell$, (ii) a token-aggregation map
$g:\{h^{(\ell)}_t\}_t \mapsto z\in\mathbb{R}^{d}$ (or $\mathbb{R}^{d'}$), and (iii) a probe
$p_\theta$ of capacity class $\mathcal{C}$ fit on $(z, y)$ pairs, where $y$ is the
gold label. The recovered predictor is $\Rc_c(M): x \mapsto p_\theta\bigl(g(\{h^{(\ell)}_t(M,x)\}_t)\bigr)$,
and the representational budget is the pair
\begin{equation}
c \;=\; (\,\text{capacity class } \mathcal{C},\ \text{aggregation } g\,),
\end{equation}
with the layer $\ell$ swept as a free coordinate. Recovery accuracy is estimated by
5-fold cross-validated accuracy of $p_\theta$, so that $\Acc(\Rc_c(M),T)$ in
Eq.~\eqref{eq:hid} is a held-out quantity, not a fit.

 
\subsection{Task, labels, and the suppression-faithful read position}
$T$ is WMDP-bio multiple choice ($N=1273$ four-way questions). The prompt is the exact
string the model was unlearned and evaluated under (lm-evaluation-harness v0.4.2, task
\texttt{wmdp\_bio}, raw \texttt{doc\_to\_text}, no chat template): a fixed preamble, the
question, the four options rendered as \texttt{A.}--\texttt{D.}, and the cue
\texttt{Answer:}. Scoring is by letter-logit: the harness compares the logits of the
single continuation tokens $\{\,\texttt{A},\texttt{B},\texttt{C},\texttt{D}\,\}$ at the
final context token (the \texttt{:} of \texttt{Answer:}). Accordingly the probe label is
the gold answer index $y\in\{0,1,2,3\}$ (chance $=0.25$), and the default read position is
that same \texttt{:} token --- the position the suppression acts on. Matching this format
is essential: a probe that reads a different position or prompt format measures a model in
a state the suppression never defended.


 
\section{Method}
 
\subsection{The probe grid}
We sweep the representational budget $c$ over the Cartesian product of two capacity classes
and four aggregation maps, each at all $33$ layers:
\begin{itemize}[itemsep=2pt]
  \item \textbf{Capacity $\mathcal{C}$:} (a) \emph{linear} --- multinomial logistic
        regression; (b) \emph{MLP} --- one hidden layer ($256$ units, weight decay
        $10^{-3}$).
  \item \textbf{Aggregation $g$:} (a) \emph{last-token} --- read the \texttt{:} position;
        (b) \emph{full-sequence mean} --- masked mean over all prompt tokens;
        (c) \emph{option-span pool} --- mean over only the \texttt{A.}--\texttt{D.} option
        tokens; (d) \emph{learned-attention} --- a single-query attention head pools over
        all positions, jointly trained with a linear or small-MLP read head.
\end{itemize}
Each aggregation tests a distinct "hiding"/obscuring hypothesis: the last-token read asks whether the
answer survives at the scored position; the mean and option-span pools ask whether it
survives at \emph{other} positions that the readout ignores; the learned-attention head
asks whether it survives at \emph{any} position the probe could choose to attend to.

\section{Results}
 
\subsection{Pipeline validity and a true negative}
The base probe clears its gate (best layer $22$, $\Acc=0.615 \gg 0.55$) with the expected
profile: chance through layer $13$, a sharp rise at layers $14$--$16$, and a plateau of
$0.59$--$0.62$ across layers $17$--$32$, matching the base model's $\sim\!0.64$ behavioral
WMDP-bio accuracy. The RMU model's behavioral accuracy is floored at $0.299$ (correct
checkpoint, suppression confirmed). A pathology audit establishes that the RMU probe's
flat-at-chance reading is a true negative and not an extraction artifact: probes converge;
RMU activation norms are healthy (within $\sim\!15\%$ of base, no collapse); RMU held-out
accuracy is indistinguishable from its shuffled-label control at every layer while base
generalizes; and base and RMU activations are bit-for-bit identical through layer $5$,
first diverging at layer $6$ --- pinning RMU's intervention to its documented steer region
(layers $5$--$7$).
 
\subsection{The grid: no recovery from RMU at any cell}
Table~\ref{tab:grid} reports the locked-layer ($\ell=22$) result for every cell. In all
eight, the base positive control passes and the RMU signal (real minus shuffle) stays
within noise of zero: at the locked layer and, equivalently, at the best layer over all
$33$. The same probes that recover $\sim\!0.6$ from base recover nothing from RMU.


\begin{table}[h]
\centering
\small
\begin{tabular}{llccc}
\toprule
Aggregation & Capacity & Base recovery & RMU signal ($@\,\ell{=}22$) & RMU best-layer signal \\
\midrule
Last-token        & linear & $0.615$ & $\approx 0$ ($\equiv$ shuffle) & --- \\
Last-token        & MLP    & $0.640$ & $\approx 0$ & --- \\
Full-seq.\ mean   & linear & $+0.105$ & $+0.011$ & $\le 0.02$ \\
Full-seq.\ mean   & MLP    & $+0.043$ & $-0.006$ & $\le 0.02$ \\
Option-span pool  & linear & $+0.120$ & $+0.028$ & $+0.045$ \\
Option-span pool  & MLP    & $+0.108$ & $+0.020$ & $+0.046$ \\
Learned-attention & linear & $+0.396$ & $-0.010$ & $+0.021$ \\
Learned-attention & MLP    & $+0.376$ & $+0.001$ & $+0.035$ \\
\bottomrule
\end{tabular}
\caption{Representational-recovery grid at the locked layer $\ell=22$. ``Base recovery''
is accuracy for the last-token cells and signal (real $-$ shuffle) for the pooled/attention
cells. RMU signal is within noise of zero everywhere; RMU's maximum signal across all
layers and all probes is $\le\!\sim\!0.05$.}
\label{tab:grid}
\end{table}


The cells are not "equally strong nulls", and we should be transparent about that. The last-token and
learned-attention cells are the \emph{strong} nulls: their base positive control is large
($0.62$ and $+0.40$ signal respectively), something like a comparably large hidden signal would have
been detected. The full-sequence-mean and option-span cells are \emph{corroborating}: mean
pooling dilutes the answer across $\sim\!150$ tokens, weakening even the base control to
$+0.10$--$+0.12$, so these cells can only rule out a hidden signal as strong as their
(weaker) positive control could detect. All four point the same way; we lead with the
strong cells.

In terms of Eq.~\eqref{eq:hid}, for every budget $c$ in the tested grid,
$\Acc(\Rc_c(\Mrmu), T)$ never reaches any threshold $\tau$ meaningfully above chance, so the
representational $\Hid_\tau(\Mrmu, T)$ is \emph{unreached across the entire tested budget
range}, whereas for the base model it is small (a linear probe at layer $22$ already clears
$\tau=0.55$).
 
\subsection{Selectivity: the null is forget-set-specific}
On the retain set MMLU-bio ($N=454$, general biology not in the forget set), the RMU model
is \emph{not} floored: behavioral accuracy is $0.676$ (base $0.714$), and the probe recovers
the answer with signal $+0.40$ at layer $22$ so essentially equal to the base signal
$+0.44$, with the two per-layer curves tracking each other almost exactly. Table~\ref{tab:sel}
juxtaposes the forget and retain results.

 \begin{table}[h]
\centering
\small
\begin{tabular}{lcc}
\toprule
Probe ($@\,\ell{=}22$) & WMDP-bio (forget) & MMLU-bio (retain) \\
\midrule
Base & $0.615$ (recovers) & $0.709$ (recovers) \\
RMU  & $0.298$ (flat, $\equiv$ shuffle) & $0.676$ (recovers, signal $+0.40$) \\
\bottomrule
\end{tabular}
\caption{Retain-selectivity control. The RMU model's representational null is specific to
the forget set; on retain biology its residual stream decodes the answer as well as base's.}
\label{tab:sel}
\end{table}
 
This rules out the two principal benign explanations of the forget-set null:
(i) RMU did not damage biological capability broadly (both behavioral and probe accuracy on
retain bio are intact); and (ii) the probe pipeline is not generically blind to RMU
activations (it recovers RMU on the retain set, failing only on the forget set).

\section{Analysis and interpretation}
 We tested Across a broad sweep of the representational-recovery axis with two capacity classes, four
token-aggregation schemes including a learned attention style one that can attend to any position, at all
$33$ layers: the WMDP-bio answer is not decodable from RMU's residual stream, while
identical probes recover it from base everywhere, and the null is selective to the forget
set. 

This seems to establish that on the representational axis  RMU's suppression is therefore \emph{deep and targeted}
at the probe budgets tested: the suppressed capability is not hidden behind a cheap read at
any tested position or moderate capacity. This at least tells us that we are not dealing with the
failure mode hypothesized for some unlearning losses, under which a probe (or a re-applied
direction) would recover the capability at near-zero cost.

\section{Limitations so far lol}
The claim is bounded to one axis and one budget range. ``No probe we tried recovers it'' is
not ``provably absent at any capacity'': the grid spans linear-to-MLP read heads and
fixed-to-learned token access, a broad but finite sweep. More importantly, the
representational axis is one of three. A capability can be representationally inaccessible
yet cheaply restorable on the \emph{optimization} axis (a few steps of adversarial
fine-tuning) or the \emph{in-context} axis (retrieval / ICL demonstrations); indeed
data-filtering work reports exactly such an axis-dependent dissociation.
\end{document}
